\printedpage{59}
% head: A. Grothendieck : Dérivateurs, Ch. XIII, Édition des Universités de Montpellier II et Paris VII

\[
\begin{tikzcd}
X \arrow[d, "p"'] & X' \arrow[l] \arrow[d] & X'' \arrow[l, "t"'] \arrow[d] \\
Y & Y' \arrow[l] & Y'' \arrow[l, "s"]
\end{tikzcd}
\quad ,
\]

if $s \in W$, then $t \in W$.

\begin{itemize}
\item[\textbullet] $W \ \subset\ \Fl(\Mo)$

A morphism $i : A \longrightarrow B$ of $\Mo$ is a $W$-\textit{cofibration} ($i \in \mathrm{Cof}_W$) if for every diagram of cocartesian squares
\[
\begin{tikzcd}
A'' \arrow[d] & A' \arrow[l, "s"'] \arrow[d] & A \arrow[l] \arrow[d, "i"] \\
B'' & B' \arrow[l, "t"] & B \arrow[l]
\end{tikzcd}
\quad ,
\]
if $s \in W$, then $t \in W$.

\item[\textbullet] $W \ \subset\ \Fl(\Mo)$

A morphism $p : X \longrightarrow Y$ of $\Mo$ is \textit{universally in} $W$, or a $W$-\textit{universal equivalence} ($p \in W_{\mathrm{univ}}$) if for every cartesian square
\[
\begin{tikzcd}
X \arrow[d, "p"'] & X' \arrow[l] \arrow[d, "p'"] \\
Y & Y' \arrow[l]
\end{tikzcd}
\quad ,
\]
the arrow $p'$ belongs to $W$.

\item[\textbullet] $W \ \subset\ \Fl(\Mo)$

A morphism $i : A \longrightarrow B$ of $\Mo$ is \textit{couniversally in} $W$, or a $W$-\textit{couniversal equivalence} ($i \in W^{\mathrm{univ}}$) if for every cocartesian square
\[
\begin{tikzcd}
A' \arrow[d, "i'"'] & A \arrow[l] \arrow[d, "i"] \\
B' & B \arrow[l]
\end{tikzcd}
\quad ,
\]
the arrow $i'$ belongs to $W$.
\end{itemize}
% foot: Version 25 janvier 2004

\printedpage{60}
% head: A. Grothendieck : Dérivateurs, Ch. XIII, Édition des Universités de Montpellier II et Paris VII

\begin{itemize}
\item[\textbullet] \textit{Left pre-Q-category of models} (resp. \textit{right}):

$\Mo$ and $W$, $C$, $F$, $TC$, $TF$ $\subset$ $\Fl(\Mo)$ satisfying axioms \textbf{M~0} to \textbf{M~5}, $\mathbf{M}'$, $\mathbf{M}''$ and \Mzg{} (resp. \Mzd).

\item[\textbullet] \textit{Q-category of models}, \textit{Quillen triple}:

model category in the sense of Quillen.

\item[\textbullet] \textit{``Closed'' Q-category of models}, \textit{closed Quillen triple}, \textit{closed Quillen category}, \textit{closed Quillen system}:

closed model category in the sense of Quillen.

\item[\textbullet] \textit{Left pre-Quillen} (resp. \textit{right}):

left pre-Q-category of models (resp. right).

\item[\textbullet] \textit{Closed Quillen pre-quintuple}:

$W$, $C$, $F$, $TC$, $TF$ $\subset$ $\Fl(\Mo)$ satisfying axioms \textbf{M~0} to \textbf{M~5}, $\mathbf{M}'$, $\mathbf{M}''$, $C$, $TC$ being left closed and $F$, $TF$ right closed.

\item[\textbullet] \textit{Weak Quillen system}:

$W$, $C$, $F$, $TC$, $TF$ $\subset$ $\Fl(\Mo)$ satisfying axioms \textbf{M~0} to \textbf{M~5}, $\mathbf{M}'$, $\mathbf{M}''$.

\item[\textbullet] \textit{Left exact weak Quillen system} (resp. \textit{right}):

$W$, $C$, $F$, $TC$, $TF$ $\subset$ $\Fl(\Mo)$ satisfying axioms \textbf{M~0} to \textbf{M~5}, $\mathbf{M}'$, $\mathbf{M}''$ and \Mzg{} (resp. \Mzd).
\end{itemize}

\section*{Section 3}

\begin{itemize}
\item[\textbullet] $(\Mo,W)$ \textit{category of models}:

$\Mo$ a category satisfying axiom \textbf{M~0};

$W$ $\subset$ $\Fl(\Mo)$ weakly saturated.

\item[\textbullet] \textit{Weak Quillen structure associated with the category of models} $(\Mo,W)$:

quadruple $(C,F,TC,TF)$ of subsets of $\Fl(\Mo)$ satisfying axioms \textbf{M~1} to \textbf{M~5}, $\mathbf{M}'$, $\mathbf{M}''$.

\item[\textbullet] \textit{Exact weak Quillen structure}:

weak Quillen structure satisfying in addition axiom \textbf{M~2\,bis}.

\item[\textbullet] \textit{Left exact weak Quillen structure} (resp. \textit{right}):

weak Quillen structure satisfying in addition axiom \Mzg{} (resp. \Mzd).
\end{itemize}
% foot: Version 25 janvier 2004

\printedpage{61}
% head: A. Grothendieck : Dérivateurs, Ch. XIII, Édition des Universités de Montpellier II et Paris VII

\begin{itemize}
\item[\textbullet] \textit{Left closed weak Quillen structure} (resp. \textit{right}):

weak Quillen structure $(C,F,TC,TF)$ such that $C$, $TC$ (resp. $F$, $TF$) are left closed (resp. right closed).

\item[\textbullet] \textit{Closed weak Quillen structure}:

weak Quillen structure both left and right closed.

\item[\textbullet] \textit{Quillen structure}:

``biexact'' weak Quillen structure:

weak Quillen structure both left and right exact

$=$ proper Quillen model category.

\item[\textbullet] \textit{Perfect Quillen structure}:

``biexact'' and closed weak Quillen structure

$=$ proper closed Quillen model category.
\end{itemize}

\section*{Section 5}

\begin{itemize}
\item[\textbullet] Category of models $(\Mo,W)$ \textit{Quillenizable}:

there exists an \textit{exact} weak Quillen structure $(C,F,TC,TF)$ associated with $(\Mo,W)$.

\item[\textbullet] $(\Mo,W)$ \textit{Quillenizable on} Diag $\subset$ Cat:

for every $I$ in Diag, the category of models $(\Mo(I),\, W(I))$ is Quillenizable.
\end{itemize}

\section*{Section 6}

\begin{itemize}
\item[\textbullet] Category of models $(\Mo,W)$ \textit{uniformly left Quillenizable on} Diag:

there exists a \textit{closed} weak Quillen structure $(C,F,TC,TF)$ for $(\Mo,\, W)$ such that for every $I$ in Diag, if one defines $C(I)$ and $TC(I)$ ``argument by argument'' (\footnote{A morphism of functors $u : X \longrightarrow Y$ from $I^{\circ}$ to $\Mo$ being in $C(I)$ (resp. $TC(I)$) if and only if for every object $i$ of $I$, $u_i$ is in $C$ (resp. $TC$).}), and $F(I)$ and $TF(I)$ by the right lifting property with respect to $TC(I)$ and $C(I)$ respectively (\footnote{A morphism of $\Mo(I)$ being in $F(I)$ (resp. $TF(I)$) if and only if it has the right lifting property with respect to the arrows belonging to $TC(I)$ (resp. $C(I)$).}), the quadruple $(C(I),TC(I),F(I),TF(I))$ is a \textit{closed exact} weak Quillen structure for $(\Mo(I),\, W(I))$. One then says that $(C,F,TC,TF)$ is a \textit{uniform left Quillenization} of $(\Mo,\, W)$ on Diag).
\end{itemize}
% foot: Version 25 janvier 2004

\printedpage{62}
% head: A. Grothendieck : Dérivateurs, Ch. XIII, Édition des Universités de Montpellier II et Paris VII

\begin{itemize}
\item[\textbullet] Category of models $(\Mo,W)$ \textit{Quillenizable on} Diag \textit{uniformly and exactly on the left}:

$(\Mo,W)$ admits a uniform left Quillenization on Diag $(C,F,TC,TF)$ such that $TC = C \cap W$ (in other words, such that $(C,F,TC,TF)$ is a \textit{left exact} weak Quillen strucure). One then says that $(C,F,TC,TF)$ is a \textit{uniform and left exact Quillenization} of $(\Mo,\, W)$ on Diag. % [sic] strucure -> structure, letter t omitted
\end{itemize}

\section*{Section 7}

\begin{itemize}
\item[\textbullet] $I$ a filtering ordered set, \quad $\pi$ a cardinal

$I$ \textit{large relative to} $\pi$:

every subset of $I$ of cardinality less than or equal to $\pi$ admits an upper bound in $I$.

\item[\textbullet] $E$ a category

A \textit{cardinal filtration} of $E$ is an increasing filtration
\[
\Big(\Filt^{\pi}(E)\Big)_{\pi\geq\pi_0} \qquad , \qquad\qquad \pi_0 \text{ a fixed infinite cardinal,}
\]
of $E$ by strictly full subcategories $\Filt^{\pi}(E)$, indexed by the cardinals $\pi \geq \pi_0$ belonging to a given universe, and satisfying the following conditions:
\begin{itemize}
\item[\textit{a})] for every $\pi \geq \pi_0$, $\Filt^{\pi}(E)$ is equivalent to a small category;
\item[\textit{b})] the category $E$ admits filtering direct limits indexed by ordered sets $I$ large relative to $\pi_0$, and for every $\pi \geq \pi_0$, $\Filt^{\pi}(E)$ is stable under the said limits if $\mathrm{card}\, I \leq \pi$;
\item[\textit{c})] For every $\pi \geq \pi_0$, and for every $X \in \mathrm{Ob}\, E$, there exists an isomorphism
\[
X \ \simeq\ \varinjlim_{i\in I} X_i \ ,
\]
where $(X_i)_{i\in I}$ is a filtering direct system in $E$, indexed by an ordered set $I$ large relative to $\pi$, and where the $X_i$ are objects of $\Filt^{\pi}(E)$. Moreover, if $X \in \mathrm{Ob}\,\Filt^{\pi'}(E)$, $\pi' \geq \pi$, one can take $I$ such that $\mathrm{card}\, I \leq \pi'^{\pi}$.
\end{itemize}

\item[\textbullet] $E$ a category, \quad $(\Filt^{\pi}(E))_{\pi\geq\pi_0}$ a cardinal filtration of $E$, \quad $\pi \geq \pi_0$ a cardinal

A \textit{well-adapted direct system} (to the given cardinal filtration, for the given cardinals $\pi_0$, $\pi$) is a filtering direct system $(X_i)_{i\in I}$ satisfying the following conditions:
\begin{itemize}
\item[\textit{a})] $I$ is large relative to $\pi$, and for every $i \in I$, $X_i \in \Filt^{\pi}(E)$;
\item[\textit{b})] Every subset $J$ of $I$ of cardinality $\leq \pi$, and large relative to $\pi_0$, admits a least upper bound $i_J$, and we have
\[
X_{i_J} \ \xleftarrow{\ \sim\ }\ \varinjlim_{i\in J} X_i \ .
\]
% foot: Version 25 janvier 2004
\end{itemize}

\printedpage{63}
% head: A. Grothendieck : Dérivateurs, Ch. XIII, Édition des Universités de Montpellier II et Paris VII

If $X$ is an object of $E$, a \textit{well-adapted representation} of $X$, in terms of objects of $\Filt^{\pi}(E)$, is the datum of a well-adapted direct system $(X_i)_{i\in I}$, to the given cardinal filtration, for the cardinals $\pi_0$, $\pi$, and of an isomorphism
\[
X \ \xleftarrow{\ \sim\ }\ \varinjlim_{i\in I} X_i \ .
\]

\item[\textbullet] $E$ a category admitting small filtering direct limits

A \textit{strict cardinal filtration} of $E$ is a cardinal filtration $(\Filt^{\pi}(E))_{\pi\geq\pi_0}$ such that for every cardinal $\pi \geq \pi_0$, $\Filt^{\pi}(E)$ is stable under filtering direct limits.

\item[\textbullet] $E$ a category admitting small filtering direct limits

$(\Filt^{\pi}(E))_{\pi\geq\pi_0}$ a strict cardinal filtration of $E$, \quad $\pi \geq \pi_0$ a cardinal

A \textit{perfectly adapted direct system} (to the given cardinal filtration, for the cardinal $\pi$) is a filtering direct system $(X_i)_{i\in I}$ satisfying the following conditions:
\begin{itemize}
\item[\textit{a})] $I$ is large relative to $\pi$, and for every $i \in I$, $X_i \in \Filt^{\pi}(E)$;
\item[\textit{b})] Every filtering subset $J$ of $I$ of cardinality $\leq \pi$, admits a least upper bound $i_J$, and we have
\[
X_{i_J} \ \xleftarrow{\ \sim\ }\ \varinjlim_{i\in J} X_i \ .
\]
\end{itemize}
If $X$ is an object of $E$, a \textit{perfectly adapted representation} of $X$, in terms of objects of $\Filt^{\pi}(E)$, is the datum of a well-adapted direct system $(X_i)_{i\in I}$, to the given cardinal filtration, for the cardinal $\pi$, and of an isomorphism
\[
X \ \xleftarrow{\ \sim\ }\ \varinjlim_{i\in I} X_i \ .
\]
\end{itemize}
% foot: Version 25 janvier 2004